
Model capacity below a performance threshold invalidates correlation-based selective prediction experiments on factual question answering---not by reducing statistical power, but by making correlation tests mathematically undefined. We demonstrate this through entropy extraction experiments on TriviaQA that succeeded at infrastructure validation (100\% extraction rate, 8.20\% high max-probability, high-entropy disagreement cases) while hypothesis testing failed (Spearman correlation = NaN due to zero-variance correctness). The root cause: GPT-2 (117M parameters) produced 0\% accuracy (0/500 correct), eliminating the variance required for correlation analysis. We tested only GPT-2; based on scaling laws literature, we hypothesize models below approximately 7B parameters fall below this validity threshold on knowledge-intensive tasks, though empirical confirmation across model scales remains future work. Our infrastructure validation succeeded independent of model performance: token probability distributions are fully extractable from frozen LLM forward passes, and entropy-max-probability disagreement patterns exist in non-trivial proportions. However, our core hypothesis---that entropy-based rejection outperforms max-probability-based rejection---remains untested due to invalid experimental conditions. We propose a methodological guideline: selective prediction experiments must verify non-zero variance in evaluation metrics before conducting correlation analyses, distinguishing undefined tests (invalid conditions) from negative tests (hypothesis unsupported). This framework applies broadly to correlation-based uncertainty quantification, where model capacity gates experiment validity rather than just affecting performance magnitude.
